% ============================================================================
% APPENDIX B: MATHEMATICAL DERIVATIONS
% ============================================================================

\chapter{Mathematical Derivations}
\label{app:math}

This appendix provides detailed mathematical derivations for key algorithms in the system.

\section{Quadrotor Dynamics}

\subsection{Coordinate Frames}

The system uses two primary coordinate frames:
\begin{itemize}
    \item \textbf{Inertial Frame (NED):} Fixed to Earth, with X pointing North, Y pointing East, Z pointing Down.
    \item \textbf{Body Frame:} Fixed to the vehicle, with X pointing forward, Y pointing right, Z pointing down.
\end{itemize}

\subsection{Rotation Matrix}

The rotation from body frame to inertial frame uses the ZYX Euler angle convention:

\begin{equation}
    \mathbf{R}_{b}^{i} = \mathbf{R}_z(\psi) \mathbf{R}_y(\theta) \mathbf{R}_x(\phi)
\end{equation}

where:
\begin{align}
    \mathbf{R}_x(\phi) &= \begin{bmatrix}
        1 & 0 & 0 \\
        0 & \cos\phi & -\sin\phi \\
        0 & \sin\phi & \cos\phi
    \end{bmatrix} \\[10pt]
    \mathbf{R}_y(\theta) &= \begin{bmatrix}
        \cos\theta & 0 & \sin\theta \\
        0 & 1 & 0 \\
        -\sin\theta & 0 & \cos\theta
    \end{bmatrix} \\[10pt]
    \mathbf{R}_z(\psi) &= \begin{bmatrix}
        \cos\psi & -\sin\psi & 0 \\
        \sin\psi & \cos\psi & 0 \\
        0 & 0 & 1
    \end{bmatrix}
\end{align}

The full rotation matrix is:
\begin{equation}
    \mathbf{R}_{b}^{i} = \begin{bmatrix}
        c_\psi c_\theta & c_\psi s_\theta s_\phi - s_\psi c_\phi & c_\psi s_\theta c_\phi + s_\psi s_\phi \\
        s_\psi c_\theta & s_\psi s_\theta s_\phi + c_\psi c_\phi & s_\psi s_\theta c_\phi - c_\psi s_\phi \\
        -s_\theta & c_\theta s_\phi & c_\theta c_\phi
    \end{bmatrix}
\end{equation}

where $c_x = \cos(x)$ and $s_x = \sin(x)$.

\subsection{Translational Dynamics}

The translational equations of motion in the inertial frame:

\begin{equation}
    m \ddot{\mathbf{p}} = \mathbf{R}_{b}^{i} \begin{bmatrix} 0 \\ 0 \\ -T \end{bmatrix} + \begin{bmatrix} 0 \\ 0 \\ mg \end{bmatrix} + \mathbf{F}_{\text{dist}}
\end{equation}

where:
\begin{itemize}
    \item $m$ is the vehicle mass
    \item $T$ is the total thrust force
    \item $g$ is gravitational acceleration
    \item $\mathbf{F}_{\text{dist}}$ is the disturbance force (e.g., wind)
\end{itemize}

Expanding:
\begin{align}
    \ddot{x} &= \frac{T}{m}(\cos\phi\sin\theta\cos\psi + \sin\phi\sin\psi) + \frac{F_{d,x}}{m} \\
    \ddot{y} &= \frac{T}{m}(\cos\phi\sin\theta\sin\psi - \sin\phi\cos\psi) + \frac{F_{d,y}}{m} \\
    \ddot{z} &= \frac{T}{m}\cos\phi\cos\theta - g + \frac{F_{d,z}}{m}
\end{align}

\subsection{Rotational Dynamics}

The angular velocity transformation from body rates to Euler rates:

\begin{equation}
    \begin{bmatrix} \dot{\phi} \\ \dot{\theta} \\ \dot{\psi} \end{bmatrix} =
    \begin{bmatrix}
        1 & \sin\phi\tan\theta & \cos\phi\tan\theta \\
        0 & \cos\phi & -\sin\phi \\
        0 & \sin\phi\sec\theta & \cos\phi\sec\theta
    \end{bmatrix}
    \begin{bmatrix} p \\ q \\ r \end{bmatrix}
\end{equation}

The rotational equations of motion in the body frame:

\begin{align}
    I_{xx}\dot{p} &= (I_{yy} - I_{zz})qr + \tau_\phi \\
    I_{yy}\dot{q} &= (I_{zz} - I_{xx})pr + \tau_\theta \\
    I_{zz}\dot{r} &= (I_{xx} - I_{yy})pq + \tau_\psi
\end{align}

\section{Extended Kalman Filter Derivation}

\subsection{State Space Model}

The continuous-time state space model:
\begin{align}
    \dot{\mathbf{x}} &= f(\mathbf{x}, \mathbf{u}) + \mathbf{w} \\
    \mathbf{z} &= h(\mathbf{x}) + \mathbf{v}
\end{align}

where $\mathbf{w} \sim \mathcal{N}(0, \mathbf{Q}_c)$ and $\mathbf{v} \sim \mathcal{N}(0, \mathbf{R})$.

\subsection{Discretization}

For discrete-time implementation with timestep $\Delta t$:

\begin{equation}
    \mathbf{x}_{k+1} = \mathbf{x}_k + f(\mathbf{x}_k, \mathbf{u}_k) \Delta t + \mathbf{w}_k
\end{equation}

The discrete process noise covariance:
\begin{equation}
    \mathbf{Q}_d = \mathbf{Q}_c \Delta t
\end{equation}

\subsection{Jacobian Computation}

The state transition Jacobian $\mathbf{F} = \frac{\partial f}{\partial \mathbf{x}}$:

\begin{equation}
    \mathbf{F} = \begin{bmatrix}
        \mathbf{0}_{3\times3} & \mathbf{I}_{3\times3} & \mathbf{0}_{3\times3} & \mathbf{0}_{3\times3} \\
        \mathbf{0}_{3\times3} & \mathbf{0}_{3\times3} & \frac{\partial \ddot{\mathbf{p}}}{\partial \boldsymbol{\Theta}} & \mathbf{0}_{3\times3} \\
        \mathbf{0}_{3\times3} & \mathbf{0}_{3\times3} & \frac{\partial \dot{\boldsymbol{\Theta}}}{\partial \boldsymbol{\Theta}} & \frac{\partial \dot{\boldsymbol{\Theta}}}{\partial \boldsymbol{\omega}} \\
        \mathbf{0}_{3\times3} & \mathbf{0}_{3\times3} & \mathbf{0}_{3\times3} & \frac{\partial \dot{\boldsymbol{\omega}}}{\partial \boldsymbol{\omega}}
    \end{bmatrix}
\end{equation}

The discrete-time Jacobian:
\begin{equation}
    \mathbf{F}_d = \mathbf{I} + \mathbf{F} \Delta t
\end{equation}

\subsection{Kalman Gain Derivation}

The Kalman gain minimizes the posterior covariance:

\begin{equation}
    \mathbf{K} = \arg\min_{\mathbf{K}} \text{tr}(\mathbf{P}^+)
\end{equation}

Taking the derivative and setting to zero:
\begin{align}
    \frac{\partial}{\partial \mathbf{K}} \text{tr}\left[(\mathbf{I} - \mathbf{K}\mathbf{H})\mathbf{P}^-(\mathbf{I} - \mathbf{K}\mathbf{H})^\top + \mathbf{K}\mathbf{R}\mathbf{K}^\top\right] &= 0 \\
    -2\mathbf{P}^-\mathbf{H}^\top + 2\mathbf{K}(\mathbf{H}\mathbf{P}^-\mathbf{H}^\top + \mathbf{R}) &= 0 \\
    \mathbf{K} &= \mathbf{P}^-\mathbf{H}^\top(\mathbf{H}\mathbf{P}^-\mathbf{H}^\top + \mathbf{R})^{-1}
\end{align}

\section{Control Barrier Function Theory}

\subsection{Forward Invariance}

\begin{theorem}[Nagumo's Theorem]
A closed set $\mathcal{S}$ is forward invariant for $\dot{\mathbf{x}} = f(\mathbf{x})$ if and only if:
\begin{equation}
    f(\mathbf{x}) \in T_{\mathcal{S}}(\mathbf{x}), \quad \forall \mathbf{x} \in \partial\mathcal{S}
\end{equation}
where $T_{\mathcal{S}}(\mathbf{x})$ is the tangent cone to $\mathcal{S}$ at $\mathbf{x}$.
\end{theorem}

\subsection{CBF Condition Derivation}

For a safe set $\mathcal{S} = \{\mathbf{x} : h(\mathbf{x}) \geq 0\}$, forward invariance requires:

\begin{equation}
    \dot{h}(\mathbf{x}) \geq 0 \quad \text{when } h(\mathbf{x}) = 0
\end{equation}

To ensure this with a margin, we require:
\begin{equation}
    \dot{h}(\mathbf{x}) \geq -\alpha(h(\mathbf{x}))
\end{equation}

where $\alpha$ is a class-$\mathcal{K}$ function.

For control-affine systems $\dot{\mathbf{x}} = f(\mathbf{x}) + g(\mathbf{x})\mathbf{u}$:
\begin{equation}
    \dot{h} = \underbrace{\frac{\partial h}{\partial \mathbf{x}} f(\mathbf{x})}_{L_f h} + \underbrace{\frac{\partial h}{\partial \mathbf{x}} g(\mathbf{x})}_{L_g h} \mathbf{u}
\end{equation}

The CBF constraint becomes:
\begin{equation}
    L_f h(\mathbf{x}) + L_g h(\mathbf{x}) \mathbf{u} \geq -\alpha(h(\mathbf{x}))
\end{equation}

\subsection{CBF-QP Derivation}

The optimization problem:
\begin{align}
    \min_{\mathbf{u}} \quad & \frac{1}{2}\|\mathbf{u} - \mathbf{u}_{\text{nom}}\|^2 \\
    \text{s.t.} \quad & L_f h_i + L_g h_i \mathbf{u} \geq -\alpha_i h_i, \quad i = 1, \ldots, m
\end{align}

This is a convex QP with linear constraints, solvable in polynomial time.

\subsection{Altitude Barrier Derivation}

For the altitude constraint $z_{\min} \leq z \leq z_{\max}$:

Upper bound barrier:
\begin{equation}
    h_1(z) = z - z_{\min}
\end{equation}

Derivative:
\begin{equation}
    \dot{h}_1 = \dot{z} = v_z
\end{equation}

CBF constraint:
\begin{equation}
    v_z \geq -\alpha_1(z - z_{\min})
\end{equation}

Lower bound barrier:
\begin{equation}
    h_2(z) = z_{\max} - z
\end{equation}

CBF constraint:
\begin{equation}
    -v_z \geq -\alpha_2(z_{\max} - z) \quad \Rightarrow \quad v_z \leq \alpha_2(z_{\max} - z)
\end{equation}

\section{H-Infinity Control Theory}

\subsection{Problem Formulation}

Given a generalized plant:
\begin{align}
    \dot{\mathbf{x}} &= \mathbf{A}\mathbf{x} + \mathbf{B}_1\mathbf{w} + \mathbf{B}_2\mathbf{u} \\
    \mathbf{z} &= \mathbf{C}_1\mathbf{x} + \mathbf{D}_{11}\mathbf{w} + \mathbf{D}_{12}\mathbf{u} \\
    \mathbf{y} &= \mathbf{C}_2\mathbf{x} + \mathbf{D}_{21}\mathbf{w} + \mathbf{D}_{22}\mathbf{u}
\end{align}

Find controller $\mathbf{K}$ such that:
\begin{equation}
    \|T_{zw}\|_\infty = \sup_\omega \bar{\sigma}(T_{zw}(j\omega)) < \gamma
\end{equation}

\subsection{Riccati Equation Solution}

The H-$\infty$ controller exists if and only if the following Riccati equations have stabilizing solutions:

\begin{align}
    \mathbf{A}^\top\mathbf{X} + \mathbf{X}\mathbf{A} + \mathbf{C}_1^\top\mathbf{C}_1 +
    \mathbf{X}(\gamma^{-2}\mathbf{B}_1\mathbf{B}_1^\top - \mathbf{B}_2\mathbf{B}_2^\top)\mathbf{X} &= 0 \\
    \mathbf{A}\mathbf{Y} + \mathbf{Y}\mathbf{A}^\top + \mathbf{B}_1\mathbf{B}_1^\top +
    \mathbf{Y}(\gamma^{-2}\mathbf{C}_1^\top\mathbf{C}_1 - \mathbf{C}_2^\top\mathbf{C}_2)\mathbf{Y} &= 0
\end{align}

with $\mathbf{X} \geq 0$, $\mathbf{Y} \geq 0$, and $\rho(\mathbf{X}\mathbf{Y}) < \gamma^2$.

\section{Contract Algebra}

\subsection{Contract Composition}

Given contracts $\mathcal{C}_1 = (\mathcal{A}_1, \mathcal{G}_1)$ and $\mathcal{C}_2 = (\mathcal{A}_2, \mathcal{G}_2)$:

\begin{equation}
    \mathcal{C}_1 \otimes \mathcal{C}_2 = \left(
        (\mathcal{A}_1 \cap \mathcal{A}_2) \cup (\overline{\mathcal{G}_1} \cap \mathcal{A}_2) \cup (\mathcal{A}_1 \cap \overline{\mathcal{G}_2}),
        \mathcal{G}_1 \cap \mathcal{G}_2
    \right)
\end{equation}

For polyhedral contracts (Pacti), this simplifies to constraint intersection and union operations.

\subsection{Contract Refinement}

Contract $\mathcal{C}_1$ refines $\mathcal{C}_2$ (written $\mathcal{C}_1 \preceq \mathcal{C}_2$) if:
\begin{equation}
    \mathcal{A}_2 \subseteq \mathcal{A}_1 \quad \text{and} \quad \mathcal{G}_1 \subseteq \mathcal{G}_2
\end{equation}

Interpretation: $\mathcal{C}_1$ has weaker assumptions and stronger guarantees.

